% =============================================================================
% Example 01 -- Minimal response letter
% The smallest useful document: title page, one reviewer, three comments.
% Author names, manuscript numbers and content are fictional; the journal
% name is a real journal used only as an example.
% Compile with pdfLaTeX (twice).
% =============================================================================
\documentclass[11pt]{authorresponseletter}
\usepackage[T1]{fontenc}
\usepackage{lmodern}
\usepackage{hyperref}

\title{A Placeholder Study of Example Systems}
\author{John Doe and Jane Doe}
\journal{PLOS ONE}
\manuscript{PONE-D-00-00001}

\begin{document}
\maketitle

\reviewer

\begin{generalcomment}
The manuscript studies an interesting placeholder problem and is generally well written. Some methodological details need clarification before it can be accepted.
\end{generalcomment}

% No optional argument: the default general response is printed.
\begin{generalresponse}
\end{generalresponse}

\begin{revcomment}
The data collection procedure is not described in sufficient detail.
\end{revcomment}
% No optional argument: an automatic thank-you phrase is printed.
\begin{revresponse}
We have expanded Section~2 with a full description of the data collection procedure, including the sampling interval and the cleaning steps.
\end{revresponse}

\begin{revcomment}
Figure~3 is difficult to read.
\end{revcomment}
% Optional argument: your own thank-you sentence.
\begin{revresponse}[We thank the reviewer for this helpful remark on readability.]
Figure~3 has been redrawn with larger labels and a colour-blind-friendly palette.
\end{revresponse}

\begin{revcomment}
There is a typo in Equation~(4).
\end{revcomment}
% Empty optional argument: no thank-you sentence.
\begin{revresponse}[]
Corrected.
\end{revresponse}

\end{document}
